\documentclass{article}
\usepackage[T1]{fontenc}
\usepackage{lmodern}
\usepackage{graphicx}
\usepackage{amsmath,amssymb}
\usepackage[a4paper,margin=2cm]{geometry}
\usepackage[absolute,overlay]{textpos}
\setlength{\TPHorizModule}{1mm}
\setlength{\TPVertModule}{1mm}
\usepackage[indonesian]{babel}
\usepackage{hyperref}
\setlength{\emergencystretch}{2em}
% unit: Halaman 34

\begin{document}

% === Halaman 34 (python/native) ===

2. OAEP (Optimal Asymmetric Encryption Padding)

• OAEP adalah metode padding yang ditambahkan ke pesan sebelum dienkripsi 

dengan RSA.

• Padding diperlukan karena:

• RSA tidak boleh mengenkripsi pesan langsung

• Tanpa padding, RSA rentan terhadap berbagai serangan kriptografi

• OAEP membuat pesan menjadi acak (randomized) sebelum dienkripsi RSA, 

sehingga aman terhadap chosen-ciphertext attack

• Skema ini diperkenalkan oleh Mihir Bellare dan Phillip Rogaway pada tahun 1994.

• Pesan M sebelum dienkripsi RSA diproses sebagai berikut:

       M → OAEP padding → encoded message → RSA encryption

• Sehingga ciphertext menjadi tidak deterministik (setiap enkripsi berbeda 

walaupun pesannya sama).

34

\end{document}
